\section{Four potential norms, one pair of inward tests}\label{sec:coefficient-tests}

We will construct four spaces satisfying the same midpoint bound. The
path constraint will keep their path vectors bounded; the quadratic
constraints will permit the same pair of inward perturbations.

Let $T=\mathbb N^{<\omega}$, with root $o$, and let $T^+=T\setminus\{o\}$.
Write $r\preceq s$ when $r$ is a predecessor of $s$, write $s^-$ for
the parent of a nonroot node $s$, and write $\operatorname{ch}(r)$ for
the children of $r$. On either $I=T$ or
$I=T^+$, the potential of a coefficient family $f$ is
\[
 P_f(s)=\sum_{u\in I,\ u\preceq s}f_u.
\]
There are four norming sets. Their common condition is
$|P_f(s)|\le1$ at every node; their other conditions are as follows:
\[
\begin{array}{c|c|c|l}
 K&I&\text{support condition}&\text{quadratic condition}\\ \hline
 K_{\rm s}&T^+&\text{none}&
       \sum_{u\in\operatorname{ch}(r)}|f_u|^2\le1\quad(r\in T)\\
 K_{\rm p}&T&\text{none}&
       \sum_{u\in\operatorname{ch}(r)}|f_u|^2\le1\quad(r\in T)\\
 K_{\rm a}&T&f\in c_{00}(T)&
       \sum_{u\in A}|f_u|^2\le1\quad(A\text{ a finite antichain})\\
 K_{\rm g}&T&f\in\ell_2(T)&
       \sum_{u\in T}|f_u|^2\le1.
\end{array}
\]
An antichain has no two comparable nodes. The root coefficient is
present in $K_{\rm p}$, even though it belongs to no sibling group.
Its absolute value is bounded by the path condition at $o$.


The root-inclusive sibling constraint also has a useful potential description.
A family $b:T\to[-1,1]$ with
$\sum_{u\in\operatorname{ch}(r)}|b_u-b_r|^2\le1$ determines exactly one
$f\in K_{\rm p}$ by $f_o=b_o$ and $f_s=b_s-b_{s^-}$ for $s\ne o$.
Conversely $b=P_f$ recovers every member of $K_{\rm p}$.
Thus the root label $b_o$ is free in $[-1,1]$; the root coefficient
records that label relative to a dummy parent whose value is fixed at
zero. Section~\ref{sec:free-root} omits this root coordinate and studies
the resulting norm on the edge coordinates. This convention differs
from $K_{\rm s}$, whose root potential is fixed at zero.

For each row, put $\|x\|_K=\sup_{f\in K}\sum_{s\in I}f_sx_s$ on
$c_{00}(I)$ and denote the completion by $X_K$. Each test set is symmetric, its coefficients have
absolute value at most one, and each signed singleton is a test. Thus
this formula defines a norm, $\|e_s\|_K=1$, and coordinate evaluations
are continuous. Each test extends to the completion with norm at most
one. Moreover these extensions remain norming: the supremum of their
evaluations is a $1$-Lipschitz function which agrees with the norm on a
dense subspace. We will write $f(x)$ for the extended evaluation.

The global constraint gives the additional comparison
$\|z\|_{K_{\rm g}}\le\|z\|_2$ for every finite array $z$.
Consequently the restriction of any $\varphi\in X_{K_{\rm g}}^*$
to the finite arrays is Hilbert-continuous. Its coordinate values
form an $\ell_2(T)$ family. Thus every sequence of distinct coordinate
vectors, and in particular every sibling sequence, is weakly null
in $X_{K_{\rm g}}$.

\begin{theorem}\label{thm:four}
Each of the four completed spaces $X_K$ is a separable infinite-dimensional
real Banach space. For $K\in\{K_{\rm s},K_{\rm p},K_{\rm a},K_{\rm g}\}$
and $0<t<1$,
\begin{equation}\label{eq:common-bound}
 \widehat\delta_{X_K}(t)\ge\sqrt{1+t^2/4}-1.
\end{equation}
None of these spaces admits an equivalent AUC norm.
\end{theorem}

The proof has two parts. We first obtain the midpoint bound by perturbing
a norming test outside a finite set of coordinates. We then construct
bounded paths whose child increments prevent AUC renorming.

A finite set $D\subset T$ is \emph{initial} if it contains the root
and every predecessor of each of its nodes. Restricting a test to $D$
preserves every displayed condition: quadratic sums decrease, while a
path sum freezes when the path leaves $D$. Consequently, if $x$ is
supported in $D\cap I$, there is a test $f$ supported there with
$f(x)=\|x\|_K$. Indeed we can take the supremum over tests on this fixed
finite set, which form a closed bounded set in a finite-dimensional
space. The tail space
\[
 F_D=\{y\in X_K:y_s=0\text{ for }s\in D\cap I\}
\]
is closed and has finite codimension. Its finite vectors are dense:
subtract the finitely many head coordinates from any finite
approximation to $y\in F_D$.

A one-generation positive-part perturbation appears in the proof of
Theorem~2.4 of Dilworth et al. \cite{DKRRZ2016}: they select a positive
or negative part of a Hilbert tail and orient the dual perturbation
inward from the root value. The next lemma clips a relative potential
on each exit component and pairs both inward parts.

\begin{lemma}[Paired inward tests]\label{thm:constrained-paired}
For any of the four norms, any $x$ supported in a finite initial $D$,
any $y\in F_D$, and $t>0$, one has
\begin{equation}\label{eq:pair-inequality}
 \frac{\|x+ty\|_K+\|x-ty\|_K}{2}
 \ge \sqrt{1-\theta^2}\,\|x\|_K
       +\frac{t\theta}{2}\|y\|_K
       \qquad(0\le\theta\le\tfrac12).
\end{equation}
\end{lemma}

\begin{proof}
Choose a head test $f$ with $f(x)=\|x\|_K$ and an arbitrary test $g$.
We construct two inward increment families from $g$. On each component of $T\setminus D$,
let $r$ be its last ancestor in $D$, put $c=P_f(r)$, and use the relative
potential
\[
 R(s)=P_g(s)-P_g(r),\qquad |R(s)|\le2.
\]
Let $\sigma=-1$ if $c\ge0$ and $\sigma=1$ if $c<0$, and write
$\xi_+=\max\{\xi,0\}$ for $\xi\in\mathbb R$. Define
\[
 U(s)=\sigma(\sigma R(s))_+,
 \qquad V(s)=\sigma(-\sigma R(s))_+,
\]
and set $U=V=0$ on $D$. Thus both $U$ and $V$ point inward from $c$,
and $U-V=R$. Let $u,v$ be their edge increments, with root
coefficient zero if the root is a coordinate.

Taking a signed positive part is $1$-Lipschitz and fixes zero. Along
an edge within a component the reference ancestor is unchanged; on
an edge leaving $D$ its parent has relative potential zero. It follows
in both cases that
\[
 |u_s|\le|g_s|,\qquad |v_s|\le|g_s|\quad(s\notin D),
 \qquad u_s-v_s=g_s\quad(s\notin D).
\]
Each quadratic constraint on $u$ or $v$ is therefore at most one,
using the same coordinate groups as in the chosen row of the table.
Their potentials are bounded by two, so $u/2,v/2\in K$. Thus $u$ and
$v$ extend to continuous functionals on $X_K$ with norm at most two.
Coefficient domination also preserves finite support for $K_{\rm a}$
and square summability for $K_{\rm g}$. The coefficient identity gives
$(u-v)(y)=g(y)$ first on finite tails and then, by continuity, on $F_D$.

Put $a=\sqrt{1-\theta^2}$. We claim $af+\theta u,af+\theta v\in K$.
On a tail component their potentials are $ac+\theta U(s)$ and
$ac+\theta V(s)$. If $c\ge0$, each lies between $ac-2\theta$ and $ac$;
if $c<0$, each lies between $ac$ and $ac+2\theta$. Both intervals lie
in $[-1,1]$ because $a\le1$ and $2\theta\le1$. On $D$ the potentials
are simply those of $af$. For any quadratic group $B$, the head and
tail supports are disjoint, so
\[
 \sum_{s\in B}|af_s+\theta u_s|^2
 =a^2\sum_{s\in B\cap D}|f_s|^2
   +\theta^2\sum_{s\in B\setminus D}|u_s|^2
 \le a^2+\theta^2=1.
\]
The same calculation applies to $v$; all support requirements remain
satisfied. This proves the claim for all four rows without identifying
any two of their norming sets.

Evaluate $x+ty$ against $af+\theta u$, and $x-ty$ against
$af+\theta v$. Since $f$ vanishes on tails and $u,v$ vanish on the head,
adding gives
\[
 \frac{\|x+ty\|_K+\|x-ty\|_K}{2}
 \ge a\|x\|_K+\frac{t\theta}{2}(u-v)(y)
 =a\|x\|_K+\frac{t\theta}{2}g(y).
\]
Take the supremum over $g\in K$ to obtain
\eqref{eq:pair-inequality}.
\end{proof}

The main gain from pairing is that neither endpoint must be chosen in
advance. A tail potential which is unfavorable for one endpoint supplies
an inward test for the other. The argument does not lose the tail norm
by first choosing one of the two clipped pieces.



\begin{proof}[Proof of the midpoint estimate]
Fix $t\in(0,1)$. For a finite unit center $x$ and any $y\in F_D$ with
$\|y\|_K\ge1$, choose
$\theta=t/\sqrt{4+t^2}<1/2$ in Lemma~\ref{thm:constrained-paired}. Its right
side becomes $\sqrt{1+t^2/4}$. Now let $x_0$ be any unit vector in the
completion. For every $\varepsilon>0$, choose a finite unit vector $x$
with $\|x-x_0\|_K<\varepsilon$, then choose $D$ containing its support.
The triangle inequality loses at most $\varepsilon$ in the average,
uniformly over all $y\in F_D$ of norm at least one. Therefore
\[
 \sup_{F\in\operatorname{cof}(X_K)}\inf_{\substack{y\in F\\\|y\|_K\ge1}}
 \left(\frac{\|x_0+ty\|_K+\|x_0-ty\|_K}{2}-1\right)
 \ge\sqrt{1+t^2/4}-1-\varepsilon.
\]
The subspace is allowed to depend on $\varepsilon$. Letting
$\varepsilon\downarrow0$ and then taking the infimum over $x_0$ proves
\eqref{eq:common-bound}.
\end{proof}

The common bound is quadratic: for $0<t<1$,
\[
 \sqrt{1+t^2/4}-1
 =\frac{t^2}{4(\sqrt{1+t^2/4}+1)}\ge\frac{t^2}{16}.
\]
It follows from the same admissibility argument for all four test sets;
no equality of their norms or optimality of their moduli is asserted.


\subsection{Why every equivalent norm still has bounded branches}

It remains to rule out every equivalent AUC norm. The next lemma
uses the one-sided modulus defined in \eqref{eq:auc-modulus}. Its mechanism
is simple: an AUC norm would force a fixed multiplicative increase at
every step along a suitably chosen path.

\begin{lemma}\label{lem:bounded-tree}
Let $X$ be a real Banach space. Suppose $0<m\le M<\infty$ and
$c>0$, and, for every integer $h\ge1$, $X$ contains a tree
$(p_s)_{s\in\mathbb N^{\le h}}$ with $m\le\|p_s\|\le M$.
For every nonterminal $s$, suppose the child increments
$d_{s,n}=p_{s^\frown n}-p_s$ are weakly null as $n\to\infty$ and
$\|d_{s,n}\|\ge c$. If $N$ is an equivalent norm with
$\alpha\|\cdot\|\le N\le\beta\|\cdot\|$ for constants
$0<\alpha\le\beta<\infty$, then
$\overline\delta_N(\alpha c/(2\beta M))=0$.
\end{lemma}

\begin{proof}
Write $t_0=\alpha c/(2\beta M)$ and suppose
$0<\gamma<\overline\delta_N(t_0)$. At any current node $p_s$, the
definition supplies a finite-codimensional $F$ such that
\[
 N(p_s+w)\ge(1+\gamma)N(p_s)
 \quad(w\in F,\ N(w)=t_0N(p_s)).
\]
The same inequality holds when $N(w)\ge t_0N(p_s)$. To see this,
write a shorter point on the ray as a convex combination of $p_s$ and
$p_s+w$; convexity gives
$N(p_s+rw)-N(p_s)\ge r[N(p_s+w)-N(p_s)]$ for $r\ge1$.

The quotient $X/F$ is finite dimensional. Since $d_{s,n}$ is weakly
null, choose $w_n\in F$ with $N(w_n-d_{s,n})\to0$. Eventually
$N(w_n)\ge\alpha c/2\ge t_0N(p_s)$. Choosing the approximation error
at most $\gamma N(p_s)/2$ gives a child satisfying
\[
 N(p_{s^\frown n})\ge(1+\gamma/2)N(p_s).
\]
Choose $h$ large enough that $\alpha m(1+\gamma/2)^h>\beta M$,
and perform the selection in a tree of height $h$. Repeating it $h$ times
yields a path whose last vector has
$N$-norm at least $\alpha m(1+\gamma/2)^h$. This exceeds the uniform bound $\beta M$, a contradiction. Finally,
$\overline\delta_N(t_0)\ge0$: at a unit center intersect any candidate
subspace with the kernel of a norming functional, on which the
perturbed norms are at least one. This proves equality to zero.
\end{proof}

\begin{proof}[Completion of Theorem~\ref{thm:four}]
The coordinate vectors are linearly independent, so each space is
infinite dimensional. Finite rational coordinate combinations are dense,
so each space is separable. For $s\in I$, set
$p_s=\sum_{u\in I,\ u\preceq s}e_u$. Its norm is one: the path
condition gives the upper bound, and a singleton test gives the lower
bound. In the rootless model begin at any fixed nonroot node $v$ and
use the subtree $(p_{v^\frown s})_s$ so that the initial vector is
nonzero. In all four models,
\[
 \left\|\sum_{n\in A}a_ne_{s^\frown n}\right\|_K
 \le\left(\sum_{n\in A}|a_n|^2\right)^{1/2}
 \quad(A\subset\mathbb N\text{ finite}).
\]
This is the sibling constraint in the first two models. Siblings form
an antichain in the third and are part of the globally constrained
family in the fourth. If $\phi\in X_K^*$, use
$a_n=\phi(e_{s^\frown n})$ in this inequality to obtain
\[
 \sum_{n\in A}|\phi(e_{s^\frown n})|^2\le\|\phi\|^2.
\]
Thus the child increments are weakly null, and each has norm one.
Lemma~\ref{lem:bounded-tree} applies with $m=M=c=1$. In particular,
for every equivalent norm as in that lemma,
$\overline\delta_N(\alpha/(2\beta))=0$.
\end{proof}
